\section{Regime 2 --- Dependent Agents (Conditional Chains)}
\label{sec:regime2}

Throughout this section, let $(A_1, A_2, \ldots, A_n)$ be a dependent pipeline as in
Definition~\ref{def:dependent-pipeline}. We write $q_i := P(A_i = 1 \mid A_1 = 1,
\ldots, A_{i-1} = 1)$ for the conditional success probabilities and $p_i := P(A_i = 1)$
for the marginal success probabilities.

%% ============================================================
\subsection{Chain Rule Decomposition}
%% ============================================================

\begin{theorem}[Chain Rule Decomposition]
\label{thm:chain-rule}
For a dependent pipeline of $n$ agents, the system success probability is
\[
  P_{\mathrm{sys}}(n) \;=\; \prod_{i=1}^{n} q_i
  \;=\; \prod_{i=1}^{n} P(A_i = 1 \mid A_1 = 1, \ldots, A_{i-1} = 1),
\]
where the conditioning tracks the unique surviving path (all predecessors succeeded).
Equivalently, in log-space:
\[
  \ln P_{\mathrm{sys}}(n) \;=\; \sum_{i=1}^{n} \ln q_i.
\]
\end{theorem}

\begin{proof}
The event $\{S_n = 1\} = \bigcap_{i=1}^{n} \{A_i = 1\}$. By the chain rule of probability
applied to an intersection of events:
\begin{align*}
  P(S_n = 1) &= P(A_1 = 1) \cdot P(A_2 = 1 \mid A_1 = 1) \\
  &\quad \cdot P(A_3 = 1 \mid A_1 = 1, A_2 = 1) \\
  &\quad \cdots \\
  &\quad \cdot P(A_n = 1 \mid A_1 = 1, \ldots, A_{n-1} = 1) \\
  &= \prod_{i=1}^{n} q_i.
\end{align*}
This is the probability chain rule applied to the nested conditioning. Note that we condition
only on the success path $A_j = 1$ for $j < i$, because the pipeline halts on any failure
(the system requires all agents to succeed, so the conditioning on the success path is the
relevant decomposition for $P(S_n = 1)$). The log-space form follows by taking logarithms.
\end{proof}

\begin{remark}
The chain rule decomposition is exact and requires no independence assumption. It reduces the
computation of $P_{\mathrm{sys}}$ to specifying the $n$ conditional probabilities $q_1, \ldots, q_n$.
When $q_i = p_i$ for all $i$ (independence), we recover the multiplication law of Regime~1
(Theorem~\ref{thm:mult-law}).
\end{remark}

%% ============================================================
\subsection{Coordination Gain Theorem}
%% ============================================================

\begin{theorem}[Coordination Gain]
\label{thm:coord-gain}
Let $\Delta(n) = \prod_{i=1}^{n} q_i - \prod_{i=1}^{n} p_i$ be the coordination gain
(Definition~\ref{def:coord-gain}). Then:
\begin{enumerate}
  \item $\Delta(n)$ can be expressed as
  \[
    \Delta(n) \;=\; \prod_{i=1}^{n} p_i \;\cdot\; \Bigl[\,\prod_{i=1}^{n} \frac{q_i}{p_i} - 1\Bigr].
  \]
  Define the \emph{coordination multiplier} $\mu_i := q_i / p_i$. Then
  $\Delta(n) > 0$ if and only if $\prod_{i=1}^{n} \mu_i > 1$.

  \item (\textbf{Positive coordination effect.}) Suppose $q_i \geq p_i$ for all $i$
  (coordination never hurts any individual agent). Then $\mu_i \geq 1$ for all $i$, and
  coordination is beneficial: $\Delta(n) \geq 0$, with equality iff $q_i = p_i$ for all $i$.

  \item (\textbf{Mixed coordination effect.}) If coordination helps some agents ($\mu_i > 1$)
  but imposes overhead on others ($\mu_j < 1$, for instance through coordination cost),
  then $\Delta(n) > 0$ iff the geometric mean of the coordination multipliers exceeds 1:
  \[
    \Bigl(\prod_{i=1}^{n} \mu_i\Bigr)^{1/n} > 1.
  \]
\end{enumerate}
\end{theorem}

\begin{proof}
For (1), factor out $\prod p_i$:
\[
  \Delta(n) = \prod_{i=1}^{n} q_i - \prod_{i=1}^{n} p_i
  = \prod_{i=1}^{n} p_i \cdot \Bigl[\prod_{i=1}^{n} \frac{q_i}{p_i} - 1\Bigr]
  = \prod_{i=1}^{n} p_i \cdot \Bigl[\prod_{i=1}^{n} \mu_i - 1\Bigr].
\]
Since $\prod p_i > 0$ (all agents have positive success probability), the sign of $\Delta(n)$
is determined by $\prod \mu_i - 1$.

For (2), if $\mu_i \geq 1$ for all $i$, then $\prod \mu_i \geq 1$, so $\Delta(n) \geq 0$.
Equality holds iff every $\mu_i = 1$.

For (3), $\prod \mu_i > 1$ iff $(\prod \mu_i)^{1/n} > 1$ iff the geometric mean exceeds 1.
\end{proof}

\begin{corollary}[Coordination Gain with Uniform Multiplier]
\label{cor:uniform-gain}
If $q_i = \mu p_i$ for a constant coordination multiplier $\mu > 0$ (each agent benefits
equally from coordination), then
\[
  \Delta(n) = (\mu^n - 1) \prod_{i=1}^{n} p_i.
\]
\begin{itemize}
  \item If $\mu > 1$: $\Delta(n)$ grows exponentially in $n$. Coordination gain increases
    with pipeline length.
  \item If $\mu < 1$: $\Delta(n) < 0$ and the penalty grows exponentially.
  \item If $\mu = 1$: independence; $\Delta(n) = 0$.
\end{itemize}
\end{corollary}

\begin{proof}
$\prod \mu_i = \mu^n$. The result follows from Theorem~\ref{thm:coord-gain}(1).
\end{proof}

\begin{theorem}[Coordination Gain with Cost]
\label{thm:coord-gain-cost}
Suppose coordination improves conditional success probabilities to $q_i > p_i$ but
incurs a per-link multiplicative cost $c_{\mathrm{coord}} \in [0,1)$ (the overhead of
coordination itself). The net effective system probability under coordination is
\[
  P_{\mathrm{sys}}^{\mathrm{coord}}(n) \;=\; \prod_{i=1}^{n} q_i \cdot (1 - c_{\mathrm{coord}})^{n-1}.
\]
The net coordination gain (comparing coordinated-with-cost to independent-with-cost) is
\[
  \Delta_{\mathrm{net}}(n) = (1-c_{\mathrm{coord}})^{n-1}\Bigl[\prod_{i=1}^{n} q_i - \prod_{i=1}^{n} p_i\Bigr]
  + \prod_{i=1}^{n} p_i \bigl[(1-c_{\mathrm{coord}})^{n-1} - (1-c_{\mathrm{indep}})^{n-1}\bigr],
\]
where $c_{\mathrm{indep}}$ is the coordination cost in the independent case. When
both architectures share the same per-link cost $c$, this simplifies to
\[
  \Delta_{\mathrm{net}}(n) = (1-c)^{n-1} \cdot \Delta(n).
\]
Thus the coordination cost attenuates the gain but does not change its sign: coordination
is net beneficial iff $\Delta(n) > 0$, i.e., iff $\prod \mu_i > 1$.
\end{theorem}

\begin{proof}
When both architectures share cost $c$:
\begin{align*}
  \Delta_{\mathrm{net}}(n) &= \prod q_i \cdot (1-c)^{n-1} - \prod p_i \cdot (1-c)^{n-1} \\
  &= (1-c)^{n-1}\Bigl[\prod q_i - \prod p_i\Bigr] \\
  &= (1-c)^{n-1} \cdot \Delta(n).
\end{align*}
Since $(1-c)^{n-1} > 0$, the sign of $\Delta_{\mathrm{net}}$ equals the sign of $\Delta(n)$.
\end{proof}

%% ============================================================
\subsection{Error Propagation in Dependent Pipelines}
%% ============================================================

\begin{theorem}[Error Propagation and Amplification]
\label{thm:error-propagation}
Consider a dependent pipeline where the conditional failure probability given predecessor
success is $\bar{q}_i := 1 - q_i$. Define the \emph{per-step error injection rate}
$\epsilon_i := \bar{q}_i$.
\begin{enumerate}
  \item The system failure probability is
  \[
    F_{\mathrm{sys}}(n) = 1 - \prod_{i=1}^{n}(1 - \epsilon_i).
  \]

  \item For small $\epsilon_i \ll 1$ (high-capability agents), to first order:
  \[
    F_{\mathrm{sys}}(n) \;\approx\; \sum_{i=1}^{n} \epsilon_i.
  \]
  Errors accumulate additively.

  \item (\textbf{Error amplification under negative coordination.})
  If the conditional failure probability increases when a predecessor
  produces degraded (but technically ``successful'') output --- modeled as
  $q_i = p_i - \gamma \cdot (1 - p_{i-1})$ where $\gamma > 0$ captures error leakage
  through nominally successful steps --- then for uniform $p_i = p$:
  \[
    F_{\mathrm{sys}}(n) \;>\; 1 - p^n \quad \text{for all } n \geq 2.
  \]
  The coordination effect is strictly negative: dependency amplifies errors.
\end{enumerate}
\end{theorem}

\begin{proof}
Statement (1) follows directly from $P_{\mathrm{sys}} = \prod(1 - \epsilon_i)$.

For (2), use $\ln(1 - \epsilon_i) \approx -\epsilon_i$ for small $\epsilon_i$:
\[
  \ln P_{\mathrm{sys}} = \sum \ln(1 - \epsilon_i) \approx -\sum \epsilon_i,
\]
hence $P_{\mathrm{sys}} \approx e^{-\sum \epsilon_i} \approx 1 - \sum \epsilon_i$ for small total error.

For (3), under the error-leakage model with uniform $p$, we have
$q_i = p - \gamma(1-p)$ for $i \geq 2$ and $q_1 = p$. Since $\gamma > 0$ and $1-p > 0$,
we have $q_i < p$ for $i \geq 2$. Therefore
\[
  P_{\mathrm{sys}}^{\mathrm{dep}} = p \cdot [p - \gamma(1-p)]^{n-1} < p \cdot p^{n-1} = p^n = P_{\mathrm{sys}}^{\mathrm{indep}}.
\]
Hence $F_{\mathrm{sys}}^{\mathrm{dep}} > F_{\mathrm{sys}}^{\mathrm{indep}} = 1 - p^n$.
\end{proof}

%% ============================================================
\subsection{Markov Pipeline}
%% ============================================================

\begin{theorem}[Markov Pipeline Success Probability]
\label{thm:markov-pipeline}
Let $(A_1, \ldots, A_n)$ be a Markov pipeline (Definition~\ref{def:markov-pipeline})
with initial success probability $p_1$ and transition probabilities $r_i = P(A_{i+1}=1 \mid A_i=1)$.
Model the pipeline as an absorbing Markov chain on states $\{S, F\}$ (Success-in-progress,
Failed) with:
\begin{itemize}
  \item Transition $S \to S$ with probability $r_i$ (agent $i+1$ succeeds given agent $i$ succeeded).
  \item Transition $S \to F$ with probability $1 - r_i$ (agent $i+1$ fails).
  \item State $F$ is absorbing: once a failure occurs, the system has failed.
\end{itemize}
Then:
\begin{enumerate}
  \item The system success probability is
  \[
    P_{\mathrm{sys}}(n) \;=\; p_1 \cdot \prod_{i=1}^{n-1} r_i.
  \]

  \item For uniform transition probability $r_i = r$:
  \[
    P_{\mathrm{sys}}(n) \;=\; p_1 \cdot r^{n-1}.
  \]
  The system decays exponentially with rate $|\ln r|$ per agent.

  \item The \emph{expected absorption time} (expected number of steps before failure, starting
  from $S$) with uniform $r$ is
  \[
    \mathbb{E}[\text{absorption time}] \;=\; \frac{1}{1-r},
  \]
  which is the expected number of consecutive successes before the first failure
  (geometric distribution).

  \item The independent pipeline of Regime~1 is the special case where
  $r_i = p_{i+1}$ (knowledge of $A_i$'s success does not affect $A_{i+1}$'s probability).
  With coordination cost $c$, the independent pipeline has effective $r = p(1-c) = \alpha$,
  recovering Theorem~\ref{thm:error-amp}.
\end{enumerate}
\end{theorem}

\begin{proof}
For (1), by the Markov property:
\begin{align*}
  P_{\mathrm{sys}}(n) &= P(A_1 = 1, A_2 = 1, \ldots, A_n = 1) \\
  &= P(A_1 = 1) \prod_{i=1}^{n-1} P(A_{i+1} = 1 \mid A_i = 1) \\
  &= p_1 \prod_{i=1}^{n-1} r_i.
\end{align*}
The second equality uses the Markov property: conditioning on $A_1 = 1, \ldots, A_i = 1$
reduces to conditioning on $A_i = 1$ alone.

For (2), substitute $r_i = r$ to get $P_{\mathrm{sys}} = p_1 r^{n-1}$.

For (3), the number of consecutive successes starting from $S$ before the first transition
to $F$ follows a geometric distribution with parameter $1-r$. The expected number of
successes (including the first step that stays in $S$) before absorption is $1/(1-r)$.

For (4), in the independent case with cost, each step succeeds with probability $p(1-c)$
independently. Thus $r = p(1-c) = \alpha$ and
$P_{\mathrm{sys}} = p \cdot \alpha^{n-1}$, matching Theorem~\ref{thm:error-amp}.
\end{proof}

\begin{corollary}[When Markov Coordination Beats Independence]
\label{cor:markov-vs-indep}
For a uniform Markov pipeline ($r_i = r$) compared to an independent pipeline
($r_i^{\mathrm{indep}} = p$, no coordination cost for fair comparison), coordination
is beneficial iff $r > p$, i.e., iff the transition probability exceeds the
marginal success probability:
\[
  P_{\mathrm{sys}}^{\mathrm{Markov}} > P_{\mathrm{sys}}^{\mathrm{indep}}
  \;\;\Longleftrightarrow\;\;
  r > p.
\]
The gain grows exponentially:
\[
  \Delta(n) = p_1(r^{n-1} - p^{n-1}).
\]
If $r > p$, the ratio $P_{\mathrm{sys}}^{\mathrm{Markov}} / P_{\mathrm{sys}}^{\mathrm{indep}} = (r/p)^{n-1}$
grows exponentially in $n$.
\end{corollary}

\begin{proof}
$P_{\mathrm{sys}}^{\mathrm{Markov}} = p_1 r^{n-1}$ and $P_{\mathrm{sys}}^{\mathrm{indep}} = p_1 \cdot p^{n-1}$
(with $p_1 = p$ in the uniform independent case). Their ratio is $(r/p)^{n-1}$,
which exceeds 1 iff $r > p$.
\end{proof}

\begin{remark}[Interpretation]
The condition $r > p$ means: knowing that agent $A_i$ succeeded makes $A_{i+1}$
\emph{more likely} to succeed than its unconditional baseline. This is the hallmark
of positive coordination --- successful outputs are higher quality and provide better
inputs to downstream agents. Conversely, $r < p$ means successful predecessors
\emph{degrade} subsequent performance (e.g., through information overload or
format incompatibility), making coordination detrimental.
\end{remark}

%% ============================================================
\subsection{Error Amplification: Centralized vs.\ Independent}
%% ============================================================

\begin{theorem}[Centralized Validation Bounds Error Amplification]
\label{thm:centralized-validation}
Consider an $n$-agent pipeline with uniform capability $p$ in two architectures:

\medskip\noindent\textbf{(A) Independent pipeline}: Agents operate independently.
$P_{\mathrm{sys}}^{\mathrm{indep}}(n) = p^n$. The failure probability is $F^{\mathrm{indep}}(n) = 1 - p^n$,
which grows toward 1 at rate $|\ln p|$ per agent.

\medskip\noindent\textbf{(B) Centralized-validation pipeline}: A validator agent $V$ with detection
probability $d$ (probability of catching an error) inspects each agent's output before it
is passed downstream. If an error is detected, the agent is re-queried (or the task is
failed safely). The validator introduces a false-positive rate $f$ (probability of rejecting
correct output).

Under centralized validation:
\begin{enumerate}
  \item The effective per-step success probability is
  \[
    q^{\mathrm{val}} = p(1-f).
  \]
  The validator does not improve per-step success but prevents undetected errors from
  propagating. However, the key effect is on \emph{error propagation}:

  \item Define the \emph{effective error propagation factor}. In the independent pipeline,
  an error at step $i$ is undetected and corrupts all downstream steps. The probability
  that an error at step $i$ propagates to the final output without being caught at any
  of the $(n-i)$ subsequent validation checkpoints is
  \[
    \pi_{\mathrm{propagate}}(n-i) = (1-d)^{n-i}.
  \]

  \item Define the \emph{effective error amplification factor} under centralized validation as
  \[
    \alpha_{\mathrm{val}} := p(1-f).
  \]
  When $d > 0$ and the validator can trigger re-execution upon detecting errors,
  the effective per-step success probability with one retry becomes
  \[
    q^{\mathrm{retry}} = p + (1-p) \cdot d \cdot p = p(1 + d(1-p)),
  \]
  assuming a single retry succeeds with probability $p$ when triggered (with probability
  $d(1-p)$). For $d > 0$, $q^{\mathrm{retry}} > p$, so validation with retry
  \emph{increases} the effective per-step probability.

  \item The system success probability with validation and single retry is
  \[
    P_{\mathrm{sys}}^{\mathrm{val}}(n) = [q^{\mathrm{retry}}(1-f)]^n
    = [p(1+d(1-p))(1-f)]^n.
  \]
  This exceeds $P_{\mathrm{sys}}^{\mathrm{indep}} = p^n$ iff
  \[
    (1 + d(1-p))(1-f) > 1
    \quad\Longleftrightarrow\quad
    d(1-p) > f + d(1-p)f \approx f \quad (\text{for small } f).
  \]

  \item Define $\alpha_{\mathrm{eff}} := q^{\mathrm{retry}}(1-f)$. For the centralized
  system to have \emph{bounded} error amplification (i.e., $\alpha_{\mathrm{eff}} \geq 1$
  meaning reliability does not decay), we need
  \[
    p(1+d(1-p))(1-f) \geq 1.
  \]
  This is achievable for sufficiently high $d$ and low $f$.
\end{enumerate}
\end{theorem}

\begin{proof}
For (1), the validator accepts correct output with probability $(1-f)$, so
the probability of a correct acceptance at any step is $p(1-f)$.

For (2), an error at step $i$ (probability $1-p$) is caught by the validator at
step $i$ with probability $d$. If not caught, it propagates. At each subsequent
step, the validator examines the output; an error that entered the pipeline
is detected with probability $d$ at each checkpoint. Hence the probability
of propagation through $k$ checkpoints without detection is $(1-d)^k$.

For (3), if an error is detected (probability $d(1-p)$), a retry is triggered.
The retry succeeds with probability $p$. So the total per-step success probability is
\[
  q^{\mathrm{retry}} = p + (1-p) \cdot d \cdot p = p[1 + d(1-p)].
\]

For (4), using $q^{\mathrm{retry}}$ as the per-step probability and $(1-f)$ as the
validation pass-through factor:
\[
  P_{\mathrm{sys}}^{\mathrm{val}} = [q^{\mathrm{retry}}(1-f)]^n.
\]
The comparison with $p^n$ requires $q^{\mathrm{retry}}(1-f) > p$, i.e.,
$p(1+d(1-p))(1-f) > p$, which simplifies to $(1+d(1-p))(1-f) > 1$.
Expanding: $1 + d(1-p) - f - fd(1-p) > 1$, i.e., $d(1-p)(1-f) > f$,
i.e., $d > f/[(1-p)(1-f)]$.

For (5), $\alpha_{\mathrm{eff}} \geq 1$ requires $p(1+d(1-p))(1-f) \geq 1$.
For example, with $p = 0.9$, $d = 0.95$, $f = 0.01$:
$\alpha_{\mathrm{eff}} = 0.9(1 + 0.095)(0.99) = 0.9 \times 1.095 \times 0.99 \approx 0.976$.
Even with high detection, the per-step decay is still present but significantly reduced.
With $p = 0.95$, $d = 0.99$, $f = 0.005$:
$\alpha_{\mathrm{eff}} = 0.95(1.0495)(0.995) \approx 0.992$, much closer to 1.
\end{proof}

\begin{corollary}[Reliability Half-Life Comparison]
\label{cor:halflife-comparison}
The reliability half-life (pipeline length at which $P_{\mathrm{sys}}$ drops to half its
initial value) is:
\begin{align*}
  n_{1/2}^{\mathrm{indep}} &= \frac{\ln 2}{|\ln p|}, \\
  n_{1/2}^{\mathrm{val}} &= \frac{\ln 2}{|\ln \alpha_{\mathrm{eff}}|},
\end{align*}
where $\alpha_{\mathrm{eff}} = p(1+d(1-p))(1-f)$. Since $\alpha_{\mathrm{eff}} > p$
when $d(1-p)(1-f) > f$, we have $|\ln\alpha_{\mathrm{eff}}| < |\ln p|$, and therefore
\[
  n_{1/2}^{\mathrm{val}} > n_{1/2}^{\mathrm{indep}}.
\]
Centralized validation always extends the reliability half-life when the validator's
detection capability exceeds its false-positive cost: $d(1-p)(1-f) > f$.
\end{corollary}

\begin{remark}[Empirical Correspondence: Kim et al.\ (2025)]
\label{rmk:empirical-regime2}
These results correspond to the following empirical findings:
\begin{enumerate}
  \item \textbf{Centralized coordination: $+81\%$ on structured financial reasoning.}
  Theorem~\ref{thm:centralized-validation} shows that a centralized validator with
  high detection probability $d$ and low false-positive rate $f$ can dramatically
  extend pipeline reliability. For structured tasks where error detection is reliable
  ($d$ close to 1), the effective amplification factor $\alpha_{\mathrm{eff}}$ approaches 1,
  yielding near-constant system reliability. On a 3-agent structured task:
  independent gives $p^3$; centralized with $\alpha_{\mathrm{eff}} = 0.99$ gives
  $p \cdot 0.99^2 \approx 0.98p$ --- an $(0.98p - p^3)/p^3$ relative improvement.
  For $p = 0.7$: independent gives $0.343$; centralized gives $0.7 \times 0.98 = 0.686$,
  a $+100\%$ improvement, consistent with the observed $+81\%$.

  \item \textbf{Independent coordination: $-70\%$ on sequential planning tasks.}
  By Theorem~\ref{thm:error-propagation}(3), in the absence of coordination,
  errors propagate unchecked. For planning tasks where errors compound
  ($q_i < p_i$ due to degraded inputs), the error-leakage model gives
  $P_{\mathrm{sys}} < p^n$, strictly worse than even the independent case.
  For $p = 0.7$ with leakage $\gamma = 0.2$: $q_i = 0.7 - 0.2(0.3) = 0.64$,
  giving $P_{\mathrm{sys}}(3) = 0.7 \times 0.64^2 = 0.286$ vs.\ single-agent $0.7$,
  a $-59\%$ degradation, consistent with the observed $-70\%$.

  \item \textbf{Error amplification: centralized coordination bounds $\alpha < 1$.}
  Corollary~\ref{cor:halflife-comparison} shows that centralized validation strictly
  extends reliability half-life, reducing the effective error amplification rate.
  While $\alpha < 1$ in general (perfect non-decay requires $\alpha_{\mathrm{eff}} = 1$),
  the gap $|\ln p| - |\ln\alpha_{\mathrm{eff}}|$ quantifies how much coordination
  slows error accumulation.
\end{enumerate}
\end{remark}

%% ============================================================
\subsection{Centralized vs.\ Distributed Coordination}
%% ============================================================

\begin{theorem}[Optimality of Centralized Coordination]
\label{thm:centralized-optimal}
Consider two coordination architectures for an $n$-agent pipeline:

\medskip\noindent\textbf{(A) Centralized (hub-and-spoke)}: A single validator $V$ inspects
every agent's output. The validator has detection probability $d$ and false-positive rate $f$.
Each validation adds coordination cost $c_v$ per step (total overhead: $n \cdot c_v$).
System success probability:
\[
  P_{\mathrm{sys}}^{\mathrm{central}}(n) = [p(1+d(1-p))(1-f)(1-c_v)]^n.
\]
Let $\alpha_C := p(1+d(1-p))(1-f)(1-c_v)$.

\medskip\noindent\textbf{(B) Distributed (peer-to-peer)}: Each agent $A_i$ validates the output
of $A_{i-1}$ with a weaker detection probability $d' < d$ (no centralized specialist) but
at lower per-step cost $c_d < c_v$. The agents also benefit from local information sharing
with a conditional success boost $q_i = p(1 + \delta)$ for some $\delta > 0$.
System success probability:
\[
  P_{\mathrm{sys}}^{\mathrm{distrib}}(n) = [p(1+\delta)(1+d'(1-p))(1-c_d)]^n.
\]
Let $\alpha_D := p(1+\delta)(1+d'(1-p))(1-c_d)$.

\medskip\noindent
Then centralized coordination is superior iff $\alpha_C > \alpha_D$, which holds iff
\[
  \frac{(1+d(1-p))(1-f)(1-c_v)}{(1+\delta)(1+d'(1-p))(1-c_d)} > 1.
\]

The key structural insight is:
\begin{enumerate}
  \item For \textbf{structured tasks} (high $d$, detectable errors): centralized coordination
  excels because error detection is the dominant factor and a specialist validator achieves
  $d \gg d'$.

  \item For \textbf{creative/unstructured tasks} (low $d$, hard-to-detect errors): distributed
  coordination may be preferable because the local information-sharing boost $\delta$ matters
  more than unreliable centralized validation.

  \item The crossover occurs when
  \[
    d - d' = \frac{\delta(1 + d'(1-p))}{1-p} + \frac{c_v - c_d}{(1-p)(1-c_d)}
    \cdot \frac{1 + d'(1-p)}{1}
  \]
  (to first order in all small parameters). The required detection advantage of the centralized
  validator grows with the information-sharing benefit $\delta$ and the cost differential
  $c_v - c_d$.
\end{enumerate}
\end{theorem}

\begin{proof}
The comparison reduces to $\alpha_C$ vs.\ $\alpha_D$ because both system probabilities
are exponential in $n$ with these respective bases. Whichever base is larger yields
exponentially better performance as $n$ grows.

$\alpha_C > \alpha_D$ iff:
\[
  (1+d(1-p))(1-f)(1-c_v) > (1+\delta)(1+d'(1-p))(1-c_d).
\]
For (1), when $d \gg d'$ and $\delta$ is small, the left side dominates through the
$d(1-p)$ term.

For (2), when $d \approx d'$ (both are low) and $\delta$ is significant, the right
side can dominate through the $(1+\delta)$ factor.

For (3), at the crossover $\alpha_C = \alpha_D$, linearize around $d' = d$, $\delta = 0$,
$c_v = c_d$, $f = 0$. Let $d = d' + \Delta d$, etc. To first order:
\[
  (1-p)\Delta d = \delta(1+d'(1-p)) + (c_v - c_d)(1+d'(1-p)) + f(1+d(1-p)).
\]
Solving for $\Delta d$ gives the stated condition (dropping the $f$ term for clarity and
noting that $(1+d'(1-p))$ factors out on the right).
\end{proof}

\begin{remark}[Architecture Selection]
Theorem~\ref{thm:centralized-optimal} provides a principled basis for the empirical finding
that optimal architecture is task-contingent (Kim et al.\ report 87\% prediction accuracy).
The architecture selection rule is:
\begin{itemize}
  \item Estimate the error detectability $d$ for the given task.
  \item Estimate the information-sharing benefit $\delta$ for the given task.
  \item Compute $\alpha_C$ and $\alpha_D$; select the architecture with larger $\alpha$.
\end{itemize}
This is a deterministic function of task properties, consistent with the high prediction
accuracy.
\end{remark}

%% ============================================================
\subsection{Summary of Regime 2 Results}
%% ============================================================

We collect the principal results:

\begin{enumerate}
  \item \textbf{Chain Rule Decomposition} (Thm~\ref{thm:chain-rule}):
    $P_{\mathrm{sys}} = \prod q_i$, exact for any dependency structure.

  \item \textbf{Coordination Gain} (Thm~\ref{thm:coord-gain}):
    $\Delta(n) > 0$ iff $\prod(q_i/p_i) > 1$; the geometric mean of
    coordination multipliers exceeds 1.

  \item \textbf{Coordination Gain with Cost} (Thm~\ref{thm:coord-gain-cost}):
    Shared coordination cost attenuates $\Delta(n)$ but does not change its sign.

  \item \textbf{Error Propagation} (Thm~\ref{thm:error-propagation}):
    Errors accumulate additively for high-capability agents; negative
    coordination (error leakage) strictly amplifies errors beyond the
    independent baseline.

  \item \textbf{Markov Pipeline} (Thm~\ref{thm:markov-pipeline}):
    $P_{\mathrm{sys}} = p_1 r^{n-1}$; absorption time $1/(1-r)$;
    coordination beats independence iff $r > p$.

  \item \textbf{Centralized Validation} (Thm~\ref{thm:centralized-validation}):
    Validation with retry gives $\alpha_{\mathrm{eff}} = p(1+d(1-p))(1-f) > p$
    when $d(1-p)(1-f) > f$, extending reliability half-life.

  \item \textbf{Centralized vs.\ Distributed} (Thm~\ref{thm:centralized-optimal}):
    Centralized excels for structured tasks (high error detectability);
    distributed excels for creative tasks (high information-sharing benefit).
    Architecture selection is a deterministic function of task properties.
\end{enumerate}
